\documentclass[11pt]{article}
\usepackage[margin=1in]{geometry}
\usepackage{amsmath,amssymb,amsthm}
\usepackage{algorithm}
\usepackage{algpseudocode}
\usepackage{booktabs}
\usepackage{hyperref}

\newtheorem{definition}{Definition}
\newtheorem{theorem}{Theorem}
\newtheorem{lemma}[theorem]{Lemma}
\newtheorem{property}{Property}

\title{\textbf{Mathematical Specification, DAG Tape Topology, and Adjoint Propagation of Reverse-Mode Automatic Differentiation (ALGO-NN-23)}}
\author{Team Scaibu \\ \small Email: \texttt{jaydeep.9664920749@gmail.com} \\ \small Architecture and Core Engine Team}
\date{\today}

\begin{document}
\maketitle

\begin{abstract}
Reverse-mode automatic differentiation (Reverse AD) evaluates the exact gradient $\nabla_\mathbf{x} f$ of a scalar objective $f: \mathbb{R}^n \to \mathbb{R}$ with computational complexity bounded by a small constant multiple of the forward pass cost, strictly independent of input dimension $n$. This document formalizes computational graph tapes, proves the Baur-Strassen Cheap Gradient Theorem, provides full pseudocode matching the reference implementation, and derives memory-compute tradeoffs.
\end{abstract}

\section{Notation and Mathematical Preliminaries}
\label{sec:notation}

Let $G = (V, E)$ denote a Directed Acyclic Graph (DAG) whose vertices $V = \{v_1, v_2, \dots, v_M\}$ represent intermediate variables arranged in topological order, with $v_M = \mathcal{L} \in \mathbb{R}$ denoting the terminal scalar objective.

\subsection{Adjoint State Mechanics}
\paragraph{Intuitive Meaning:} Reverse-mode AD computes gradients in reverse topological order. Beginning with seed $\bar{v}_M = \frac{\partial \mathcal{L}}{\partial \mathcal{L}} = 1.0$, each node pulls incoming derivative contributions from all of its downstream consumers.

\paragraph{Formal Mathematical Definition:}
\begin{definition}[Adjoint State]
The adjoint variable $\bar{v}_i$ for node $v_i$ is:
\begin{equation}
\label{eq:adjoint_def}
\bar{v}_i = \frac{\partial \mathcal{L}}{\partial v_i} = \sum_{j \in \operatorname{Children}(v_i)} \bar{v}_j \frac{\partial v_j}{\partial v_i}.
\end{equation}
\end{definition}

\paragraph{Worked Numerical Example:}
Let $v_1 = 2.0, v_2 = 3.0$, and $v_3 = v_1 \cdot v_2 = 6.0$. $\mathcal{L} = v_3$.
Forward: $v_3 = 6.0$.
Backward: $\bar{v}_3 = 1.0$.
$\bar{v}_1 = \bar{v}_3 \cdot \frac{\partial v_3}{\partial v_1} = 1.0 \times v_2 = 3.0$.
$\bar{v}_2 = \bar{v}_3 \cdot \frac{\partial v_3}{\partial v_2} = 1.0 \times v_1 = 2.0$.

\subsection{Summary of Symbols}
\begin{itemize}
    \item $M \in \mathbb{N}$: Total number of nodes in the computational graph tape.
    \item $\mathbf{v} \in \mathbb{R}^M$: Vector of forward evaluated scalar values.
    \item $\bar{\mathbf{v}} \in \mathbb{R}^M$: Vector of accumulated adjoint gradient values.
    \item $\operatorname{Parents}(v_j) \subset \{1, \dots, j-1\}$: Direct antecedent nodes in DAG.
\end{itemize}

\section{Problem Definition and Core Concepts}
\label{sec:problem_definition}

\subsection{Elementary Operator Jacobians}
\begin{itemize}
    \item \textbf{Addition ($v = a + b$):} $\bar{a} \gets \bar{a} + \bar{v}, \quad \bar{b} \gets \bar{b} + \bar{v}$.
    \item \textbf{Multiplication ($v = a \cdot b$):} $\bar{a} \gets \bar{a} + \bar{v} \cdot b, \quad \bar{b} \gets \bar{b} + \bar{v} \cdot a$.
    \item \textbf{ReLU ($v = \max(0, a)$):} $\bar{a} \gets \bar{a} + \bar{v} \cdot \mathbb{I}_{\{a > 0\}}$.
    \item \textbf{Sine ($v = \sin(a)$):} $\bar{a} \gets \bar{a} + \bar{v} \cdot \cos(a)$.
    \item \textbf{Exponential ($v = \exp(a)$):} $\bar{a} \gets \bar{a} + \bar{v} \cdot v$.
\end{itemize}

\section{Algorithmic Specification}
\label{sec:algorithm}

Algorithm~\ref{alg:reverse_ad} specifies the tape evaluation routine matching \texttt{impl.py}.

\begin{algorithm}[H]
\caption{Reverse-Mode Automatic Differentiation on Computational DAG Tapes}
\label{alg:reverse_ad}
\begin{algorithmic}[1]
\Require Computational tape of nodes $\mathcal{T} = [N_1, \dots, N_M]$, target loss node ID $T_{\mathrm{target}} \in [0, M-1]$.
\Ensure Forward values $\mathbf{v} \in \mathbb{R}^M$, adjoints $\bar{\mathbf{v}} \in \mathbb{R}^M$.
\State Initialize $\mathbf{v} \gets [0.0]_{i=1}^M, \quad \bar{\mathbf{v}} \gets [0.0]_{i=1}^M$
\For{$i = 0$ to $M-1$} \Comment{Forward Topological Pass}
    \State $\mathrm{op} \gets N_i.\mathrm{op}, \quad \mathrm{parents} \gets N_i.\mathrm{parents}$
    \If{$\mathrm{op} = \text{input}$}
        \State $v_i \gets N_i.\mathrm{value}$
    \ElsIf{$\mathrm{op} = \text{add}$}
        \State $v_i \gets v_{\mathrm{parents}[0]} + v_{\mathrm{parents}[1]}$
    \ElsIf{$\mathrm{op} = \text{mul}$}
        \State $v_i \gets v_{\mathrm{parents}[0]} \cdot v_{\mathrm{parents}[1]}$
    \ElsIf{$\mathrm{op} = \text{relu}$}
        \State $v_i \gets \max(0.0, v_{\mathrm{parents}[0]})$
    \ElsIf{$\mathrm{op} = \text{sin}$}
        \State $v_i \gets \sin(v_{\mathrm{parents}[0]})$
    \ElsIf{$\mathrm{op} = \text{exp}$}
        \State $v_i \gets \exp(v_{\mathrm{parents}[0]})$
    \EndIf
\EndFor
\State $\bar{v}_{T_{\mathrm{target}}} \gets 1.0$ \Comment{Reverse Adjoint Seed}
\For{$i = M-1$ down to $0$} \Comment{Reverse Adjoint Accumulation}
    \State $\bar{v} \gets \bar{v}_i$
    \If{$\bar{v} \neq 0.0$}
        \State $\mathrm{op} \gets N_i.\mathrm{op}, \quad \mathrm{parents} \gets N_i.\mathrm{parents}$
        \If{$\mathrm{op} = \text{add}$}
            \State $\bar{v}_{\mathrm{parents}[0]} \gets \bar{v}_{\mathrm{parents}[0]} + \bar{v}, \quad \bar{v}_{\mathrm{parents}[1]} \gets \bar{v}_{\mathrm{parents}[1]} + \bar{v}$
        \ElsIf{$\mathrm{op} = \text{mul}$}
            \State $\bar{v}_{\mathrm{parents}[0]} \gets \bar{v}_{\mathrm{parents}[0]} + \bar{v} \cdot v_{\mathrm{parents}[1]}, \quad \bar{v}_{\mathrm{parents}[1]} \gets \bar{v}_{\mathrm{parents}[1]} + \bar{v} \cdot v_{\mathrm{parents}[0]}$
        \ElsIf{$\mathrm{op} = \text{relu}$}
            \State $\bar{v}_{\mathrm{parents}[0]} \gets \bar{v}_{\mathrm{parents}[0]} + \bar{v} \cdot \mathbb{I}_{\{v_{\mathrm{parents}[0]} > 0\}}$
        \ElsIf{$\mathrm{op} = \text{sin}$}
            \State $\bar{v}_{\mathrm{parents}[0]} \gets \bar{v}_{\mathrm{parents}[0]} + \bar{v} \cdot \cos(v_{\mathrm{parents}[0]})$
        \ElsIf{$\mathrm{op} = \text{exp}$}
            \State $\bar{v}_{\mathrm{parents}[0]} \gets \bar{v}_{\mathrm{parents}[0]} + \bar{v} \cdot v_i$
        \EndIf
    \EndIf
\EndFor
\State \Return $\mathbf{v}, \quad \bar{\mathbf{v}}$
\end{algorithmic}
\end{algorithm}

\section{Theoretical Guarantees and Complexity Bound}
\label{sec:guarantees}

\begin{theorem}[Baur-Strassen Cheap Gradient Theorem]
\label{thm:baur_strassen_full}
\textbf{Assumptions:} Let $f: \mathbb{R}^n \to \mathbb{R}$ be evaluated via a computational graph containing $M$ primitive binary/unary operations with execution time $\operatorname{Time}(f) = \mathcal{O}(M)$.
\textbf{Guarantee:} Reverse-mode automatic differentiation computes the complete gradient vector $\nabla f(\mathbf{x}) \in \mathbb{R}^n$ with time $\operatorname{Time}(\nabla f) \le 5 \cdot \operatorname{Time}(f)$, completely independent of parameter dimension $n$.
\textbf{Proof:}
Every elementary operation $v_j = \phi(v_a, v_b)$ has at most 2 parents. Backward vector-Jacobian accumulation evaluates $\bar{v}_a \gets \bar{v}_a + \bar{v}_j \frac{\partial \phi}{\partial v_a}$ and $\bar{v}_b \gets \bar{v}_b + \bar{v}_j \frac{\partial \phi}{\partial v_b}$, requiring at most $c \le 4$ operations per node. Summing over all $M$ nodes in reverse topological order yields $\sum_{j=1}^M \operatorname{Cost}(\text{VJP}_j) \le 4M \le 4 \operatorname{Time}(f)$. Adding the initial forward pass gives $\operatorname{Time}(\nabla f) \le 5 \operatorname{Time}(f)$.
\textbf{Limitations:} Peak memory scales linearly $\mathcal{O}(M)$ to retain intermediate forward activations until the reverse pass completes.
\end{theorem}

\section{Complexity Analysis}
\label{sec:complexity}

\begin{itemize}
    \item \textbf{Time Complexity:} $\mathcal{O}(M)$ operations forward, $\mathcal{O}(M)$ operations backward. Total time is $\mathcal{O}(M)$.
    \item \textbf{Space Complexity:} $\mathcal{O}(M)$ memory for node values and adjoint storage.
\end{itemize}

\section{Worked Numerical Example}
\label{sec:worked_example}

Let graph compute $f(x_1, x_2) = x_1 \cdot x_2 + \sin(x_1)$ with $x_1 = \pi/2 \approx 1.570796, x_2 = 2.0$:
\begin{enumerate}
    \item Forward:
    $v_0 = 1.570796, v_1 = 2.0$.
    $v_2 = v_0 \cdot v_1 = 3.141592$.
    $v_3 = \sin(v_0) = 1.0$.
    $v_4 = v_2 + v_3 = 4.141592$.
    \item Backward with $\bar{v}_4 = 1.0$:
    $\bar{v}_2 = 1.0, \bar{v}_3 = 1.0$.
    From $v_3$: $\bar{v}_0 \gets 1.0 \cdot \cos(1.570796) = 0.0$.
    From $v_2$: $\bar{v}_0 \gets \bar{v}_0 + 1.0(v_1) = 0.0 + 2.0 = 2.0$.
    From $v_2$: $\bar{v}_1 \gets 1.0(v_0) = 1.570796$.
\end{enumerate}

\section{Numerical Considerations and Edge Cases}
\label{sec:numerical}

\begin{itemize}
    \item \textbf{Topological Ordering Requirement:} The execution tape must be strictly topologically sorted prior to evaluation so that every child node is visited before any of its parent nodes during the backward pass.
\end{itemize}

\begin{thebibliography}{9}
\bibitem{griewank2008} A.~Griewank and A.~Walther, \emph{Evaluating Derivatives: Principles and Techniques of Algorithmic Differentiation}, 2nd~ed., SIAM, 2008.
\end{thebibliography}

\end{document}
